% Generated draft: requires mathematical and source review.
\documentclass[11pt]{article}
\usepackage[margin=25mm]{geometry}
\usepackage{fontspec}
\setmainfont{Latin Modern Roman}
\usepackage{amsmath,amssymb,mathtools,unicode-math}
\setmathfont{Latin Modern Math}
\usepackage{enumitem,needspace,xurl,hyperref}
\setlength{\parindent}{0pt}
\setlength{\parskip}{5pt}
\setlength{\emergencystretch}{3em}
\newcommand{\N}{\mathbb{N}}
\newcommand{\Z}{\mathbb{Z}}
\newcommand{\Q}{\mathbb{Q}}
\newcommand{\R}{\mathbb{R}}
\newcommand{\C}{\mathbb{C}}
\newcommand{\F}{\mathbb{F}}
\newcommand{\A}{\mathbb{A}}
\newcommand{\PP}{\mathbb{P}}
\newcommand{\E}{\mathbb{E}}
\newcommand{\eps}{\varepsilon}
\newcommand{\dd}{\,\mathrm{d}}
\begin{document}
\section*{UnsolvedMath ID: 30001011}
\textbf{Source problem code:} catalogue entry
\subsection{Definitions and mathematical statement}
Let \(M\) be a \(\mathrm{II}_1\) factor with normalized faithful normal trace \(\tau\), so \(\tau(1)=1\). For a bounded self-adjoint element \(a\in M\), its spectral distribution is the compactly supported probability measure \(\mu_a\) on \(\R\) defined by
\[
\mu_a(B)=\tau(1_B(a))\qquad(B\subset\R\text{ Borel}),
\]
using spectral projections. Equivalent spectral data may be expressed by a decreasing quantile, or eigenvalue, function; equality of such functions is understood almost everywhere.

For this given factor \(M\), and prescribed distributions \(\mu,\nu\), characterize
\[
\{\mu_{a+b}:a=a^*,\ b=b^*\in M,\ \mu_a=\mu,\ \mu_b=\nu\}.
\]
Here the operators vary over all pairs having the prescribed marginal spectral distributions; they are not required to commute or be free. Equivalently, characterize all triples \((\mu_a,\mu_b,\mu_c)\) realized by self-adjoint \(a,b,c\in M\) satisfying \(a+b+c=0\). In passing between the two forms, \(c=-(a+b)\), so \(\mu_c\) is the pushforward of \(\mu_{a+b}\) under \(x\mapsto-x\).

The requested characterization is for an arbitrary \(\mathrm{II}_1\) factor as stated. Embeddability is not a hypothesis of the question, and no specific complete system of Horn-type inequalities or sufficiency theorem is supplied in the extraction.
\subsection{Short English statement}
For two prescribed self-adjoint spectral distributions in a fixed \(\mathrm{II}_1\) factor, determine all distributions of their sums. The zero-sum triple version uses reflection of the third distribution and imposes no commutativity or freeness.
\subsection{Sources}
\begin{itemize}
\item Local UnsolvedMath registry, permanent ID \(30001011\), source code \texttt{catalogue entry}, ``Horn Spectral Data in Finite \(\mathrm{II}_1\) Factors.'' Registry curation: Ulam AI / UnsolvedMath Contributors, CC BY 4.0; underlying sources retain their supplied provenance.
\item \emph{\(C^*\)-Algebras} (2009), Oberwolfach Reports, pp.\ 2095--2164, DOI \url{https://doi.org/10.4171/owr/2008/37}. The supplied title has malformed escaping, normalized here solely for display.
\item Supplied ``Bercovici, Collins, Dykema, Li and Timotin'': \url{https://arxiv.org/abs/0805.4817}.
\item Supplied ``Collins and Dykema'': \url{https://arxiv.org/abs/0711.3930}.
\item Supplied literature triage, \texttt{literature\_checked\_at}: 2026-08-21: labeled partially solved; says Horn-type inequalities hold for arbitrary finite factors, a full characterization is known for embeddable factors, and general sufficiency remains open and is tied to the embedding problem. This is supplied dated curation, not independently verified; its embedding-related status claims require separate review.
\end{itemize}
\end{document}
